\documentclass[a4paper]{article}
% Español (México) con XeLaTeX: fontspec para fuentes Unicode, polyglossia
% para la división silábica y los nombres del documento en la variante mexicana.
\usepackage{fontspec}
\usepackage{polyglossia}
\setdefaultlanguage[variant=mexican]{spanish}
% Los nombres chinos van como «pinyin (original)»: xeCJK da a los caracteres
% Han su propia fuente; sin ella XeLaTeX los omite con un aviso.
% FandolSong no cubre algunos caracteres de nombres (U+73FA); WenQuanYi Zen Hei sí.
\usepackage[AutoFallBack=true]{xeCJK}
\setCJKmainfont{FandolSong-Regular.otf}
\setCJKfallbackfamilyfont{\CJKrmdefault}{WenQuanYi Zen Hei}
% polyglossia no traduce los nombres de listings.
\AtBeginDocument{\providecommand{\lstlistingname}{}\renewcommand{\lstlistingname}{Listado}%
  \providecommand{\lstlistlistingname}{}\renewcommand{\lstlistlistingname}{Índice de listados}}
\usepackage{amsmath, amssymb}
\usepackage{graphicx}
\usepackage[margin=2.5cm]{geometry}
\usepackage[most]{tcolorbox}
\usepackage{etoolbox}
\usepackage{listings}
\usepackage{hyperref}
\usepackage{booktabs}
\usepackage{float}
\newtcolorbox{knowledgebox}[1]{enhanced, colback=blue!5!white, colframe=blue!75!black, colbacktitle=blue!75!black, coltitle=white, fonttitle=\bfseries, title={#1}, attach boxed title to top left={yshift=-2mm, xshift=2mm}, boxrule=1pt, sharp corners}
\newtcolorbox{importantbox}[1]{enhanced, colback=yellow!10!white, colframe=yellow!80!black, colbacktitle=yellow!80!black, coltitle=black, fonttitle=\bfseries, title={#1}, sharp corners}
\newtcolorbox{warningbox}[1]{enhanced, colback=red!5!white, colframe=red!75!black, colbacktitle=red!75!black, coltitle=white, fonttitle=\bfseries, title={#1}, sharp corners}
\lstset{language=Python, basicstyle=\ttfamily\small, keywordstyle=\color{blue}, stringstyle=\color{red!60!black}, commentstyle=\color{green!60!black}, breaklines=true, frame=single, numbers=left, numberstyle=\tiny\color{gray}, captionpos=b, extendedchars=true}
\newcommand{\notetitle}{Aplicación del algoritmo REINFORCE en modelos de lenguaje}
\newcommand{\noteauthors}{Elaboradas a partir de materiales públicos del curso}
\newcommand{\notedate}{2025}
\newcommand{\videochannel}{Wudaokou Nash (五道口纳什)}
\newcommand{\videopublishdate}{2025}
\newcommand{\videoduration}{aproximadamente 20 minutos}
\newcommand{\videourl}{}
\newcommand{\videocoverpath}{cover.jpg}
\newcommand{\repourl}{https://github.com/hqhq1025/ai-course-notes}

% Footer with repo info
\usepackage{fancyhdr}
\pagestyle{fancy}
\fancyhf{}
\fancyfoot[C]{\thepage}
\fancyfoot[R]{\footnotesize\color{gray}\href{\repourl}{AI Course Notes}}
\renewcommand{\headrulewidth}{0pt}
\renewcommand{\footrulewidth}{0pt}

\begin{document}
\begin{titlepage}\centering
{\Large Notas del curso\par}\vspace{1.2cm}{\huge\bfseries \notetitle\par}\vspace{0.8cm}{\large \noteauthors\par}\vspace{0.3cm}{\large \notedate\par}\vspace{1.2cm}
\ifdefempty{\videocoverpath}{}{\includegraphics[width=0.82\textwidth,height=0.45\textheight,keepaspectratio]{\videocoverpath}\par}
\vfill
\begin{tcolorbox}[width=0.9\textwidth, colback=black!2!white, colframe=black!60, sharp corners]
\textbf{Autor o canal del video}: \videochannel\par\textbf{Fecha de publicación}: \videopublishdate\par\textbf{Duración del video}: \videoduration\par
\ifdefempty{\videourl}{}{\textbf{Enlace del video}: \href{\videourl}{\nolinkurl{\videourl}}\par}\par
\textbf{Página del proyecto}: \href{\repourl}{\nolinkurl{\repourl}}\par
\end{tcolorbox}
\end{titlepage}
\tableofcontents\newpage

\section{Introducción}
En esta sesión se aplica el algoritmo REINFORCE a los modelos de lenguaje para lograr un efecto sencillo: sin importar qué prompt se introduzca, la salida del modelo siempre contiene una semántica específica (como ``mode).

\section{Diseño del entorno de Reward}
\subsection{Reward basado en Cosine Similarity}
\begin{importantbox}{Definición de Reward}
\begin{enumerate}
  \item El rollout del modelo de lenguaje genera una secuencia (oración)
  \item Se calcula el embedding de esa oración
  \item Se calcula el embedding de la palabra objetivo (por ejemplo, ``mode)
  \item La cosine similarity entre ambas se usa como reward
\end{enumerate}
\end{importantbox}

\section{Implementación de REINFORCE}
\subsection{Flujo del algoritmo}
\begin{enumerate}
  \item Muestrear secuencias generadas a partir de la política actual $\pi_\theta$
  \item Calcular el reward de cada secuencia (cosine similarity)
  \item Calcular el PG Loss y retropropagar para actualizar $\theta$
\end{enumerate}

\subsection{Resumen de la sección}
REINFORCE es el algoritmo PG más simple, que puede aplicarse directamente a modelos de lenguaje, lo que lleva los conceptos de RL a tareas reales de generación de texto.

\section{Nota de ingeniería: el diseño de reward debe servir a la retroalimentación del entrenamiento}
Aunque esta clase se centra en el diseño del entorno de reward y en la implementación de REINFORCE, el problema más importante en la práctica es: si el reward puede distinguir de manera estable entre conductas «que vale la pena reforzar» y «que no vale la pena reforzar». Si el reward es demasiado escaso, aparece demasiado tarde, o está desalineado con el objetivo final, incluso una implementación correcta del gradiente de política aprenderá muy lentamente.

\begin{warningbox}{Errores comunes en el diseño del entorno}
\begin{itemize}
  \item Solo existe una recompensa por el éxito final, sin retroalimentación intermedia, lo que vuelve demasiado difícil el credit assignment
  \item La recompensa es demasiado densa pero no restringe la dirección correcta, lo que permite que el modelo aproveche vacíos
  \item El estado del entorno es inestable, lo que hace que la misma acción reciba retroalimentación contradictoria en distintas rondas
\end{itemize}
\end{warningbox}

\subsection{Resumen de la sección}
El diseño del entorno y del reward determina la calidad de la señal de RL. En vez de ajustar el algoritmo sin más, conviene primero confirmar si la recompensa realmente está impulsando la aparición de la conducta objetivo.

\newpage
\section{Síntesis y ampliación}
\begin{enumerate}
  \item REINFORCE puede usarse directamente en el entrenamiento de RL de modelos de lenguaje
  \item El entorno de recompensa puede definirse con embedding similarity
  \item Esta es la base para entender PPO/GRPO, que son más complejos
\end{enumerate}
\end{document}
